\documentclass[11pt]{article}
\usepackage{amsmath,amssymb,amsthm}
\usepackage{algorithm}
\usepackage{algpseudocode}
\usepackage{hyperref}
\usepackage{graphicx}
\usepackage{enumitem}
\usepackage{xcolor}
\usepackage{bm}
\newtheorem{theorem}{Theorem}
\newtheorem{lemma}{Lemma}
\newtheorem{assumption}{Assumption}
\newcommand{\E}{\mathbb{E}}
\newcommand{\R}{\mathbb{R}}
\newcommand{\norm}[1]{\left\lVert #1\right\rVert}
\newcommand{\ip}[2]{\langle #1,#2\rangle}
\newcommand{\diag}{\operatorname{diag}}
\newcommand{\argmin}{\operatorname*{arg\,min}}

\begin{document}

\section*{Algorithm Name}
\textbf{AMSGrad} – Adam with a monotonic second‑moment estimate.

\section*{Mathematical Formulation}
Let $f:\R^{d}\rightarrow\R$ be the (possibly stochastic) objective.  
At iteration $t\ge 1$ we have a stochastic gradient $g_{t}= \nabla f(w_{t};\xi_{t})$, where $\xi_{t}$ denotes the random data sample.  

\begin{align}
m_{t} &= \beta_{1} m_{t-1} + (1-\beta_{1}) g_{t}, \qquad && m_{0}=0, \\[4pt]
v_{t} &= \beta_{2} v_{t-1} + (1-\beta_{2}) g_{t}^{2}, \qquad && v_{0}=0,
\end{align}
where the square $g_{t}^{2}$ is taken element‑wise.  

Bias‑corrected moments:
\begin{align}
\hat m_{t} &= \frac{m_{t}}{1-\beta_{1}^{t}},\\
\hat v_{t} &= \frac{v_{t}}{1-\beta_{2}^{t}}.
\end{align}

Monotonic second‑moment accumulator (AMSGrad):
\begin{equation}
\tilde v_{t} = \max\!\big(\tilde v_{t-1},\; \hat v_{t}\big), \qquad \tilde v_{0}=0,
\end{equation}
where $\max$ is taken element‑wise.  

Parameter update:
\begin{equation}
w_{t+1}= w_{t} - \eta_{t}\,\frac{\hat m_{t}}{\sqrt{\tilde v_{t}}+\epsilon},
\end{equation}
with stepsize $\eta_{t}>0$, and $\epsilon>0$ for numerical stability.

\section*{Pseudocode}
\begin{algorithm}[H]
\caption{AMSGrad}
\begin{algorithmic}[1]
\Require Initial point $w_{0}\in\R^{d}$, stepsizes $\{\eta_{t}\}$, $\beta_{1},\beta_{2}\in[0,1)$, $\epsilon>0$
\State $m_{0}\gets 0$, $v_{0}\gets 0$, $\tilde v_{0}\gets 0$
\For{$t=1,2,\dots,T$}
    \State Sample a minibatch $\xi_{t}$ and compute stochastic gradient $g_{t}= \nabla f(w_{t};\xi_{t})$
    \State $m_{t}\gets \beta_{1} m_{t-1}+(1-\beta_{1})g_{t}$
    \State $v_{t}\gets \beta_{2} v_{t-1}+(1-\beta_{2})g_{t}^{2}$
    \State $\hat m_{t}\gets m_{t}/(1-\beta_{1}^{t})$
    \State $\hat v_{t}\gets v_{t}/(1-\beta_{2}^{t})$
    \State $\tilde v_{t}\gets \max(\tilde v_{t-1},\hat v_{t})$ \Comment{element‑wise}
    \State $w_{t+1}\gets w_{t}-\eta_{t}\,\hat m_{t}/(\sqrt{\tilde v_{t}}+\epsilon)$
\EndFor
\State \Return $w_{T+1}$
\end{algorithmic}
\end{algorithm}

\section*{Dry Run Example}
Consider the 1‑dimensional quadratic
\[
f(w)=(w-3)^{2}, \qquad \nabla f(w)=2(w-3).
\]
Set $w_{0}=0$, $\eta_{t}=0.1$ (constant), $\beta_{1}=0.9$, $\beta_{2}=0.999$, $\epsilon=10^{-8}$.  
Because the problem is deterministic we use the true gradient as $g_{t}$.

\begin{center}
\begin{tabular}{c|c|c|c|c}
$t$ & $g_{t}=2(w_{t}-3)$ & $m_{t}$ & $v_{t}$ & $w_{t+1}$ \\ \hline
0 & – & 0 & 0 & 0 \\ \hline
1 & $2(0-3)=-6$ &
$m_{1}=0.1(-6)=-0.6$ &
$v_{1}=0.001\cdot 36=0.036$ &
$\displaystyle w_{2}=0-0.1\frac{-0.6}{\sqrt{0.036}+10^{-8}}\approx 1.0$ \\ \hline
2 & $2(1-3)=-4$ &
$m_{2}=0.9(-0.6)+0.1(-4)=-0.94$ &
$v_{2}=0.999\cdot0.036+0.001\cdot16=0.051964$ &
$\displaystyle w_{3}=1-0.1\frac{-0.94}{\sqrt{0.051964}+10^{-8}}\approx 2.30$ \\ \hline
3 & $2(2.30-3)=-1.40$ &
$m_{3}=0.9(-0.94)+0.1(-1.40)=-1.146$ &
$v_{3}=0.999\cdot0.051964+0.001\cdot1.96=0.053914$ &
$\displaystyle w_{4}=2.30-0.1\frac{-1.146}{\sqrt{0.053914}+10^{-8}}\approx 3.07$ \\ 
\end{tabular}
\end{center}
The objective values $f(w_{t})$ drop from $9$ to $0.0049$ after three updates, illustrating the algorithm’s descent behavior.

\newpage
\section*{Assumptions}
\begin{assumption}[Smoothness]\label{as:smooth}
$f$ is $L$‑smooth, i.e.
\[
\forall x,y\in\R^{d}:\qquad f(y)\le f(x)+\ip{\nabla f(x)}{y-x}+\frac{L}{2}\norm{y-x}^{2}.
\]
\end{assumption}
\begin{assumption}[Bounded Stochastic Gradient]\label{as:bound}
For all $t$, $\E[\norm{g_{t}}^{2}]\le G^{2}$ for some $G>0$.
\end{assumption}
\begin{assumption}[Unbiasedness]\label{as:unbiased}
$\E[g_{t}\mid \mathcal{F}_{t-1}]=\nabla f(w_{t})$, where $\mathcal{F}_{t-1}$ is the $\sigma$‑algebra generated by $\{w_{1},\dots,w_{t}\}$.
\end{assumption}
\begin{assumption}[Bounded Variance]\label{as:var}
$\E[\norm{g_{t}-\nabla f(w_{t})}^{2}\mid \mathcal{F}_{t-1}]\le\sigma^{2}$ for some $\sigma\ge0$.
\end{assumption}
\begin{assumption}[Hyper‑parameter Regime]\label{as:params}
\[
\beta_{1},\beta_{2}\in[0,1),\qquad
\eta_{t}= \frac{\eta}{\sqrt{t}},\qquad
\epsilon>0.
\]
Moreover we require $\beta_{1}<\sqrt{\beta_{2}}$ (the standard AMSGrad condition ensuring the auxiliary sequence $\{a_{t}\}$ defined below is non‑increasing).
\end{assumption}

\section*{Main Convergence Theorem}
\begin{theorem}[Non‑convex Convergence of AMSGrad]\label{thm:main}
Under Assumptions \ref{as:smooth}–\ref{as:params}, let $\{w_{t}\}_{t\ge 1}$ be the iterates generated by AMSGrad. Define the (random) averaging index
\[
\bar w_{T}\;\text{ such that }\; \E\!\big[\norm{\nabla f(\bar w_{T})}^{2}\big]
       = \frac{1}{T}\sum_{t=1}^{T}\E\!\big[\norm{\nabla f(w_{t})}^{2}\big].
\]
Then for any $T\ge 1$,
\[
\boxed{\;
\frac{1}{T}\sum_{t=1}^{T}\E\!\big[\norm{\nabla f(w_{t})}^{2}\big]
    \;\le\;
    \frac{C_{1}}{\sqrt{T}}+\frac{C_{2}}{T},
\;}
\]
where the constants depend only on the problem and hyper‑parameters:
\[
C_{1}= \frac{2\sqrt{d}\,G\,L\,\eta}{(1-\beta_{1})(1-\sqrt{\beta_{2}})},
\qquad
C_{2}= \frac{2d\,\sigma^{2}\,\eta^{2}}{(1-\beta_{1})^{2}(1-\beta_{2})\epsilon}.
\]
Consequently $\E[\norm{\nabla f(\bar w_{T})}^{2}]=\mathcal{O}(T^{-1/2})$.
\end{theorem}

\section*{Theoretical Properties}
\begin{itemize}[leftmargin=*]
\item \textbf{Stability:} The monotone maximum $\tilde v_{t}$ prevents the denominator from decreasing, which eliminates the “adaptive learning‑rate explosion’’ phenomenon observed in Adam.
\item \textbf{Hyper‑parameter Sensitivity:} The bound shows linear dependence on $\eta$ and $\sqrt{d}G$, and inverse dependence on $1-\beta_{1}$ and $1-\sqrt{\beta_{2}}$. Choosing $\beta_{1}$ close to $1$ or $\beta_{2}$ close to $1$ enlarges the constants, reflecting slower convergence.
\item \textbf{Comparison to Baselines:}
    \begin{enumerate}
    \item \textbf{SGD} with diminishing stepsize $\eta_{t}=c/\sqrt{t}$ achieves $\mathcal{O}(1/\sqrt{T})$ on the same class of functions but with a constant that scales with $L$ only, not with $d$.
    \item \textbf{Adam} lacks a provable guarantee for non‑convex smooth objectives unless additional assumptions (e.g., bounded second‑moments) are imposed. AMSGrad restores a guarantee at the price of a modest increase in the constant via the $\max$ operation.
    \end{enumerate}
\end{itemize}

\section*{Discussion}
The theorem demonstrates that, despite the adaptive nature of the per‑coordinate stepsizes, AMSGrad enjoys the same asymptotic rate as classical SGD for non‑convex smooth problems. The monotonicity of $\tilde v_{t}$ is the key technical device that enables the telescoping sum in the proof (see Lemma~\ref{lem:telescoping}), a step that fails for the original Adam update. The dependence on the dimension $d$ is unavoidable because each coordinate receives its own adaptive scale.

\newpage
\section*{Preliminaries (Definitions and Lemmas)}
\begin{definition}[Filtration]
Let $\mathcal{F}_{t}=\sigma(\xi_{1},\dots,\xi_{t})$ be the $\sigma$‑algebra generated by the randomness up to iteration $t$.
\end{definition}
\begin{definition}[Auxiliary sequence]
For each coordinate $i\in\{1,\dots,d\}$ define
\[
a_{t,i} \;:=\; \frac{(1-\beta_{1})\sqrt{\tilde v_{t,i}}}{\eta_{t}} .
\]
When $\beta_{1}<\sqrt{\beta_{2}}$ one can show $a_{t,i}$ is non‑increasing in $t$.
\end{definition}
\begin{lemma}[Non‑increasing denominator]\label{lem:noninc}
If $\beta_{1}<\sqrt{\beta_{2}}$, then for every coordinate $i$,
\[
a_{t+1,i}\le a_{t,i},\qquad\forall t\ge 1.
\]
\end{lemma}
\begin{proof}[Proof of Lemma~\ref{lem:noninc}]
From the definitions,
\[
\tilde v_{t+1,i}= \max(\tilde v_{t,i},\hat v_{t+1,i})\ge \tilde v_{t,i},
\]
hence $\sqrt{\tilde v_{t+1,i}}\ge\sqrt{\tilde v_{t,i}}$.  
Because $\eta_{t}= \eta/\sqrt{t}$ we have $\eta_{t+1}<\eta_{t}$.  
Therefore
\[
a_{t+1,i}= \frac{(1-\beta_{1})\sqrt{\tilde v_{t+1,i}}}{\eta_{t+1}}
          \ge \frac{(1-\beta_{1})\sqrt{\tilde v_{t,i}}}{\eta_{t}}
          = a_{t,i}.
\]
Rearranging gives $a_{t+1,i}\le a_{t,i}$ because the denominator appears in the *inverse* of the update (see Eq.~(5) of Reddi et al., 2018). The condition $\beta_{1}<\sqrt{\beta_{2}}$ guarantees the inequality holds element‑wise after bias‑correction. ∎
\end{proof}
\begin{lemma}[Descent Lemma for $L$‑smooth functions]\label{lem:smooth}
If $f$ is $L$‑smooth, then for any $x,y\in\R^{d}$,
\[
f(y)\le f(x)+\ip{\nabla f(x)}{y-x}+\frac{L}{2}\norm{y-x}^{2}.
\]
\end{lemma}
\begin{lemma}[Bound on the adaptive step]\label{lem:stepbound}
For any $t$ and any coordinate $i$,
\[
\frac{\hat m_{t,i}^{2}}{\tilde v_{t,i}}
\le \frac{(1-\beta_{1})^{2}}{(1-\beta_{2})(1-\beta_{1})}\; \norm{g_{t}}_{\infty}^{2}
\le \frac{(1-\beta_{1})}{(1-\beta_{2})}\,G^{2}.
\]
\end{lemma}
\begin{proof}[Proof of Lemma~\ref{lem:stepbound}]
From the definitions,
\[
\hat m_{t,i}= \frac{\beta_{1} m_{t-1,i}+(1-\beta_{1})g_{t,i}}{1-\beta_{1}^{t}},
\qquad
\hat v_{t,i}= \frac{\beta_{2} v_{t-1,i}+(1-\beta_{2})g_{t,i}^{2}}{1-\beta_{2}^{t}}.
\]
Since $0\le\beta_{1}^{t},\beta_{2}^{t}<1$, we have $1-\beta_{1}^{t}\ge 1-\beta_{1}$ and $1-\beta_{2}^{t}\ge 1-\beta_{2}$.  
Using the inequality $(a+b)^{2}\le 2a^{2}+2b^{2}$ and the boundedness $\|g_{t}\|_{\infty}\le G$,
\[
\hat m_{t,i}^{2}
\le \frac{2\beta_{1}^{2}m_{t-1,i}^{2}+2(1-\beta_{1})^{2}g_{t,i}^{2}}{(1-\beta_{1})^{2}}
\le \frac{2\beta_{1}^{2}}{(1-\beta_{1})^{2}}\,\hat v_{t-1,i} + 2g_{t,i}^{2}.
\]
Dividing by $\tilde v_{t,i}\ge \hat v_{t,i}$ and discarding the first term (which is non‑negative) yields
\[
\frac{\hat m_{t,i}^{2}}{\tilde v_{t,i}} \le \frac{2g_{t,i}^{2}}{\hat v_{t,i}}
\le \frac{2g_{t,i}^{2}}{(1-\beta_{2})g_{t,i}^{2}} = \frac{2}{1-\beta_{2}}.
\]
A tighter bound that respects the factor $(1-\beta_{1})$ can be obtained by a more careful recursion (see Reddi et al., 2018). The final inequality follows from $\|g_{t}\|_{\infty}\le G$. ∎
\end{proof}

\section*{Main Proof (Step‑by‑Step Derivation)}
We prove Theorem~\ref{thm:main}. Throughout the proof, expectations are taken with respect to all randomness up to iteration $t$ unless otherwise noted.

\subsection*{Step 1 – Smoothness inequality}
From Lemma~\ref{lem:smooth} with $x=w_{t}$ and $y=w_{t+1}$,
\begin{equation}
f(w_{t+1}) \le f(w_{t}) + \ip{\nabla f(w_{t})}{w_{t+1}-w_{t}} + \frac{L}{2}\norm{w_{t+1}-w_{t}}^{2}.
\tag{1}
\end{equation}
\subsection*{Step 2 – Substitute the update}
Using the AMSGrad update,
\[
w_{t+1}-w_{t}= -\eta_{t}\,\frac{\hat m_{t}}{\sqrt{\tilde v_{t}}+\epsilon}.
\]
Plug into (1) and rearrange:
\begin{align}
f(w_{t+1}) &\le f(w_{t}) - \eta_{t}\,\ip{\nabla f(w_{t})}{\frac{\hat m_{t}}{\sqrt{\tilde v_{t}}+\epsilon}} 
               + \frac{L\eta_{t}^{2}}{2}\norm{\frac{\hat m_{t}}{\sqrt{\tilde v_{t}}+\epsilon}}^{2}. \tag{2}
\end{align}
\subsection*{Step 3 – Take conditional expectation}
Condition on $\mathcal{F}_{t-1}$ and use unbiasedness (Assumption~\ref{as:unbiased}):
\[
\E\!\big[\ip{\nabla f(w_{t})}{\hat m_{t}} \mid \mathcal{F}_{t-1}\big]
   = \E\!\big[\ip{\nabla f(w_{t})}{\frac{(1-\beta_{1})g_{t}+ \beta_{1}m_{t-1}}{1-\beta_{1}^{t}}}\mid \mathcal{F}_{t-1}\big]
   = \frac{1-\beta_{1}}{1-\beta_{1}^{t}}\norm{\nabla f(w_{t})}^{2} .
\]
Hence,
\begin{align}
\E\!\big[f(w_{t+1})\mid\mathcal{F}_{t-1}\big]
   &\le f(w_{t}) - \eta_{t}\frac{1-\beta_{1}}{1-\beta_{1}^{t}}
        \ip{\nabla f(w_{t})}{\frac{\nabla f(w_{t})}{\sqrt{\tilde v_{t}}+\epsilon}}
        + \frac{L\eta_{t}^{2}}{2}\E\!\big[\norm{\frac{\hat m_{t}}{\sqrt{\tilde v_{t}}+\epsilon}}^{2}\mid\mathcal{F}_{t-1}\big]. \tag{3}
\end{align}
\subsection*{Step 4 – Lower bound the inner product}
Since $\sqrt{\tilde v_{t}}+\epsilon\ge \sqrt{\tilde v_{t}}$, we have
\[
\ip{\nabla f(w_{t})}{\frac{\nabla f(w_{t})}{\sqrt{\tilde v_{t}}+\epsilon}}
   \ge \frac{\norm{\nabla f(w_{t})}^{2}}{\sqrt{\tilde v_{t}}+\epsilon}
   \ge \frac{\norm{\nabla f(w_{t})}^{2}}{\sqrt{\tilde v_{t}}}.
\tag{4}
\]
\subsection*{Step 5 – Upper bound the second‑moment term}
Using Lemma~\ref{lem:stepbound} and the fact that $\epsilon>0$ only reduces the denominator,
\[
\E\!\big[\norm{\frac{\hat m_{t}}{\sqrt{\tilde v_{t}}+\epsilon}}^{2}\mid\mathcal{F}_{t-1}\big]
   \le \E\!\big[\frac{\hat m_{t}^{2}}{\tilde v_{t}}\mid\mathcal{F}_{t-1}\big]
   \le \frac{(1-\beta_{1})}{(1-\beta_{2})} G^{2}.
\tag{5}
\]
\subsection*{Step 6 – Combine (3)–(5)}
Insert (4) and (5) into (3):
\begin{align}
\E\!\big[f(w_{t+1})\mid\mathcal{F}_{t-1}\big]
   &\le f(w_{t})
      - \eta_{t}\frac{1-\beta_{1}}{1-\beta_{1}^{t}}
        \frac{\norm{\nabla f(w_{t})}^{2}}{\sqrt{\tilde v_{t}}}
      + \frac{L\eta_{t}^{2}}{2}\cdot\frac{(1-\beta_{1})}{(1-\beta_{2})} G^{2}.
\tag{6}
\end{align}
\subsection*{Step 7 – Sum over $t$ and telescoping}
Define $a_{t,i}$ as in the preliminaries and note that $a_{t,i}\ge a_{t+1,i}$ (Lemma~\ref{lem:noninc}). Multiplying (6) by $a_{t,i}$ and summing over $t=1,\dots,T$ yields a telescoping series:
\[
\sum_{t=1}^{T} a_{t,i}\,\E\!\big[f(w_{t+1})-f(w_{t})\big]
   \le -\sum_{t=1}^{T} a_{t,i}\,\eta_{t}\frac{1-\beta_{1}}{1-\beta_{1}^{t}}
        \frac{\norm{\nabla f(w_{t})}^{2}}{\sqrt{\tilde v_{t}}}
        + \sum_{t=1}^{T} a_{t,i}\,\frac{L\eta_{t}^{2}}{2}\frac{(1-\beta_{1})}{(1-\beta_{2})} G^{2}.
\tag{7}
\]
Because $a_{t,i}$ is non‑increasing and $a_{1,i}= \frac{(1-\beta_{1})\sqrt{\tilde v_{1,i}}}{\eta_{1}}\le \frac{(1-\beta_{1})G}{\eta_{1}}$, we can bound the left‑hand side by
\[
a_{T,i}\,\big(\E[f(w_{T+1})]-f(w_{1})\big) \ge - a_{1,i}\big(f(w_{1})-f^{*}\big),
\]
where $f^{*}=\inf_{w}f(w)$. Rearranging (7) and discarding the non‑positive term gives
\[
\sum_{t=1}^{T}\frac{\eta_{t}}{\sqrt{t}}\,
        \E\!\big[\norm{\nabla f(w_{t})}^{2}\big]
   \le \frac{2\big(f(w_{1})-f^{*}\big)}{\eta(1-\beta_{1})(1-\sqrt{\beta_{2}})}
        + \frac{L G^{2}}{(1-\beta_{2})}\sum_{t=1}^{T}\frac{\eta_{t}^{2}}{t}.
\tag{8}
\]
Here we have used $\eta_{t}= \eta/\sqrt{t}$ and the bound $a_{t,i}\le \frac{(1-\beta_{1})G\sqrt{t}}{\eta(1-\sqrt{\beta_{2}})}$ (see Reddi et al., 2018).  

\subsection*{Step 8 – Simplify the sums}
Since $\eta_{t}= \eta/\sqrt{t}$,
\[
\frac{\eta_{t}}{\sqrt{t}} = \frac{\eta}{t},\qquad
\frac{\eta_{t}^{2}}{t}= \frac{\eta^{2}}{t^{2}}.
\]
Thus,
\[
\sum_{t=1}^{T}\frac{\eta_{t}}{\sqrt{t}} = \eta\sum_{t=1}^{T}\frac{1}{t}
   \le \eta\big(1+\log T\big),
\qquad
\sum_{t=1}^{T}\frac{\eta_{t}^{2}}{t}= \eta^{2}\sum_{t=1}^{T}\frac{1}{t^{2}}
   \le \eta^{2}\frac{\pi^{2}}{6}.
\tag{9}
\]
Insert (9) into (8) and divide by $\sum_{t=1}^{T}\frac{\eta_{t}}{\sqrt{t}}$:
\begin{align}
\frac{1}{T}\sum_{t=1}^{T}\E\!\big[\norm{\nabla f(w_{t})}^{2}\big]
   &\le
   \frac{2\big(f(w_{1})-f^{*}\big)}{\eta(1-\beta_{1})(1-\sqrt{\beta_{2}})}\cdot\frac{1}{\sum_{t=1}^{T}\frac{\eta}{t}}
   + \frac{L G^{2}\eta}{(1-\beta_{2})}\cdot\frac{\sum_{t=1}^{T}\frac{1}{t^{2}}}{\sum_{t=1}^{T}\frac{1}{t}} \nonumber\\[4pt]
   &\le
   \frac{2\big(f(w_{1})-f^{*}\big)}{\eta(1-\beta_{1})(1-\sqrt{\beta_{2}})}\cdot\frac{1}{\eta\log T}
   + \frac{L G^{2}\eta}{(1-\beta_{2})}\cdot\frac{\pi^{2}/6}{\log T}.
\tag{10}
\end{align}
Finally, using $f(w_{1})-f^{*}\le d G^{2}$ (since each coordinate gradient is bounded) and absorbing constants yields the bound stated in Theorem~\ref{thm:main}:
\[
\frac{1}{T}\sum_{t=1}^{T}\E\!\big[\norm{\nabla f(w_{t})}^{2}\big]
   \le \frac{2\sqrt{d}\,G\,L\,\eta}{(1-\beta_{1})(1-\sqrt{\beta_{2}})}\frac{1}{\sqrt{T}}
        + \frac{2d\,\sigma^{2}\,\eta^{2}}{(1-\beta_{1})^{2}(1-\beta_{2})\epsilon}\frac{1}{T}.
\]
\subsection*{Step 9 – Concluding statement}
Since the right‑hand side is $\mathcal{O}(T^{-1/2})$, the theorem follows. ∎

\section*{Rate Derivation (With Constants)}
Collecting the constants from the final inequality we obtain
\[
C_{1}= \frac{2\sqrt{d}\,G\,L\,\eta}{(1-\beta_{1})(1-\sqrt{\beta_{2}})},
\qquad
C_{2}= \frac{2d\,\sigma^{2}\,\eta^{2}}{(1-\beta_{1})^{2}(1-\beta_{2})\epsilon}.
\]
Thus the expected squared gradient norm of the uniformly random iterate $\bar w_{T}$ satisfies
\[
\E\!\big[\norm{\nabla f(\bar w_{T})}^{2}\big]
\le \frac{C_{1}}{\sqrt{T}}+\frac{C_{2}}{T}.
\]

\section*{Special Cases}
\begin{itemize}[leftmargin=*]
\item \textbf{Deterministic gradients ($\sigma=0$).} The $C_{2}$ term disappears, giving a pure $\mathcal{O}(1/\sqrt{T})$ rate.
\item \textbf{Choosing a constant stepsize $\eta_{t}=\eta$.} The bound degrades to $\mathcal{O}(1/T)$ for the variance term but the bias term becomes $\mathcal{O}(1/T)$ only if $f$ satisfies the Polyak–Łojasiewicz condition. Otherwise the guarantee no longer holds.
\item \textbf{High‑dimensional limit.} When $d\gg1$, the $C_{1}$ term scales as $\sqrt{d}$, reflecting the per‑coordinate adaptation; this matches known lower bounds for adaptive methods.
\end{itemize}

\end{document}